% Part: first-order-logic
% Chapter: sequent-calculus
% Section: structural-rules

\documentclass[../../../include/open-logic-section]{subfiles}

\begin{document}

\iftag{FOL}
      {\olfileid{fol}{seq}{srl}}
      {\olfileid{pl}{seq}{srl}}

\olsection{నిర్మాణాత్మక నియమాలు}

సీక్వెంట్ యొక్క ఎడమ, కుడి వైపుల్లోని !!{sentence}sను
పునర్వ్యవస్థీకరించనిచ్చే కొన్ని నియమాలు కూడా మనకు కావాలి. తార్కిక
నియమం పనిచేసే పూర్వాధార !!{sentence}s అన్నిటికంటే ఎడమ లేదా
అన్నిటికంటే కుడివైపు ఉండాలని ఆ నియమాలు కోరతాయి. అందువల్ల
!!{sentence}sను తగిన స్థానానికి కదిలించడానికి “మార్పిడి” నియమం
అవసరం. ఒకేలా ఉన్న రెండు !!{sentence}sను ఒక్కటిగా కలపగలగడం,
ఏ వైపునైనా ఒక !!a{sentence}ను చేర్చగలగడం కూడా కొన్నిసార్లు ముఖ్యం.

\subsection{బలహీనీకరణ}

\begin{defish}
\Axiom$ \Gamma \fCenter \Delta $
\RightLabel{\LeftR{\Weakening}}
\UnaryInf$ !A, \Gamma \fCenter \Delta$
\DisplayProof
\hfill
\Axiom$ \Gamma \fCenter \Delta$
\RightLabel{\RightR{\Weakening}}
\UnaryInf$ \Gamma \fCenter \Delta, !A$
\DisplayProof
\end{defish}

\subsection{సంకోచనం}

\begin{defish}
\Axiom$ !A, !A, \Gamma \fCenter \Delta $
\RightLabel{\LeftR{\Contraction}}
\UnaryInf$ !A, \Gamma \fCenter \Delta$
\DisplayProof
\hfill
\Axiom$ \Gamma \fCenter \Delta, !A, !A$
\RightLabel{\RightR{\Contraction}}
\UnaryInf$ \Gamma \fCenter \Delta, !A$
\DisplayProof
\end{defish}

\subsection{మార్పిడి}

\begin{defish}
\Axiom$ \Gamma, !A, !B, \Pi \fCenter \Delta $
\RightLabel{\LeftR{\Exchange}}
\UnaryInf$ \Gamma, !B, !A, \Pi \fCenter \Delta$
\DisplayProof
\hfill
\Axiom$ \Gamma \fCenter \Delta, !A, !B, \Lambda$
\RightLabel{\RightR{\Exchange}}
\UnaryInf$ \Gamma \fCenter \Delta, !B, !A, \Lambda$
\DisplayProof
\end{defish}

బలహీనీకరణ, సంకోచనం, మార్పిడి నిగమనాల వరుసను తరచుగా ద్విరేఖ
నిగమన గీతలతో సూచిస్తాం.

“కట్” అని పిలిచే కింది నియమం కచ్చితంగా అవసరమైనది కాకపోయినా,
!!{derivation}sను తిరిగి ఉపయోగించడాన్ని, కలపడాన్ని ఎంతో సులభం చేస్తుంది.
\begin{defish}
\[
\Axiom$ \Gamma \fCenter \Delta, !A$
\Axiom$ !A, \Pi \fCenter \Lambda $
\RightLabel{\Cut}
\BinaryInf$ \Gamma, \Pi \fCenter \Delta, \Lambda$
\DisplayProof
\]
\end{defish}

\end{document}
